\documentclass[../../main.tex]{subfiles}
\begin{document}

{\bf [解答]}

% TODO: 座標系の概要図

$O$を原点とし，$B$を$x$軸正方向，$C$を$y$軸正方向とする座標系を考える．$P(0,(1+k)t)$とおく．
すると，まず$0\le(1+k)t\le1$より
\begin{align}
  0\le t\le\dfrac{1}{1+k}\label{eq:1}
\end{align}
である．また，内分点の条件から$Q(0,t)$であるから，$S$の$x$座標は
\begin{align}
  x^2+t^2=1 \\
  x=\sqrt{1-t^2}
\end{align}
である．故に四辺形の面積$f(t)$として
\begin{align}
  f(t)&=\dfrac{1}{2}|OP||RS| \\
  &=|OP||QS| \\
  &=(1+k)t\sqrt{1-t^2} \\
  &=(1+k)\sqrt{-\left(t^2-\dfrac{1}{2}\right)^2+\dfrac{1}{4}}
\end{align}
グラフの概形は下図のようである．

% TODO: 作図
%  \begin{center}
%    \scalebox{.5}{%WinTpicVersion4.32a
{\unitlength 0.1in%
\begin{picture}(20.0000,24.0000)(4.0000,-28.0000)%
% FUNC 2 0 3 0 Black White  
% 9 400 400 2400 2800 800 2400 2000 2400 800 1200 400 400 2400 2800 0 2 0 0
% -5x^2+5x
\special{pn 8}%
\special{pa 725 2800}%
\special{pa 725 2798}%
\special{pa 730 2770}%
\special{pa 735 2743}%
\special{pa 740 2715}%
\special{pa 745 2688}%
\special{pa 750 2660}%
\special{pa 755 2633}%
\special{pa 760 2607}%
\special{pa 765 2580}%
\special{pa 790 2450}%
\special{pa 810 2350}%
\special{pa 835 2230}%
\special{pa 840 2207}%
\special{pa 845 2183}%
\special{pa 850 2160}%
\special{pa 855 2138}%
\special{pa 860 2115}%
\special{pa 865 2093}%
\special{pa 870 2070}%
\special{pa 875 2048}%
\special{pa 880 2027}%
\special{pa 885 2005}%
\special{pa 910 1900}%
\special{pa 930 1820}%
\special{pa 955 1725}%
\special{pa 960 1707}%
\special{pa 965 1688}%
\special{pa 970 1670}%
\special{pa 975 1653}%
\special{pa 980 1635}%
\special{pa 985 1618}%
\special{pa 990 1600}%
\special{pa 995 1583}%
\special{pa 1000 1567}%
\special{pa 1005 1550}%
\special{pa 1030 1470}%
\special{pa 1050 1410}%
\special{pa 1075 1340}%
\special{pa 1080 1327}%
\special{pa 1085 1313}%
\special{pa 1090 1300}%
\special{pa 1095 1288}%
\special{pa 1100 1275}%
\special{pa 1105 1263}%
\special{pa 1110 1250}%
\special{pa 1115 1238}%
\special{pa 1120 1227}%
\special{pa 1125 1215}%
\special{pa 1150 1160}%
\special{pa 1170 1120}%
\special{pa 1195 1075}%
\special{pa 1200 1067}%
\special{pa 1205 1058}%
\special{pa 1210 1050}%
\special{pa 1215 1043}%
\special{pa 1220 1035}%
\special{pa 1225 1028}%
\special{pa 1230 1020}%
\special{pa 1235 1013}%
\special{pa 1240 1007}%
\special{pa 1245 1000}%
\special{pa 1270 970}%
\special{pa 1290 950}%
\special{pa 1315 930}%
\special{pa 1320 927}%
\special{pa 1325 923}%
\special{pa 1330 920}%
\special{pa 1335 918}%
\special{pa 1340 915}%
\special{pa 1345 913}%
\special{pa 1350 910}%
\special{pa 1355 908}%
\special{pa 1360 907}%
\special{pa 1365 905}%
\special{pa 1390 900}%
\special{pa 1410 900}%
\special{pa 1435 905}%
\special{pa 1440 907}%
\special{pa 1445 908}%
\special{pa 1450 910}%
\special{pa 1455 913}%
\special{pa 1460 915}%
\special{pa 1465 918}%
\special{pa 1470 920}%
\special{pa 1475 923}%
\special{pa 1480 927}%
\special{pa 1485 930}%
\special{pa 1510 950}%
\special{pa 1530 970}%
\special{pa 1555 1000}%
\special{pa 1560 1007}%
\special{pa 1565 1013}%
\special{pa 1570 1020}%
\special{pa 1575 1028}%
\special{pa 1580 1035}%
\special{pa 1585 1043}%
\special{pa 1590 1050}%
\special{pa 1595 1058}%
\special{pa 1600 1067}%
\special{pa 1605 1075}%
\special{pa 1630 1120}%
\special{pa 1650 1160}%
\special{pa 1675 1215}%
\special{pa 1680 1227}%
\special{pa 1685 1238}%
\special{pa 1690 1250}%
\special{pa 1695 1263}%
\special{pa 1700 1275}%
\special{pa 1705 1288}%
\special{pa 1710 1300}%
\special{pa 1715 1313}%
\special{pa 1720 1327}%
\special{pa 1725 1340}%
\special{pa 1750 1410}%
\special{pa 1770 1470}%
\special{pa 1795 1550}%
\special{pa 1800 1567}%
\special{pa 1805 1583}%
\special{pa 1810 1600}%
\special{pa 1815 1618}%
\special{pa 1820 1635}%
\special{pa 1825 1653}%
\special{pa 1830 1670}%
\special{pa 1835 1688}%
\special{pa 1840 1707}%
\special{pa 1845 1725}%
\special{pa 1870 1820}%
\special{pa 1890 1900}%
\special{pa 1915 2005}%
\special{pa 1920 2027}%
\special{pa 1925 2048}%
\special{pa 1930 2070}%
\special{pa 1935 2093}%
\special{pa 1940 2115}%
\special{pa 1945 2138}%
\special{pa 1950 2160}%
\special{pa 1955 2183}%
\special{pa 1960 2207}%
\special{pa 1965 2230}%
\special{pa 1990 2350}%
\special{pa 2010 2450}%
\special{pa 2035 2580}%
\special{pa 2040 2607}%
\special{pa 2045 2633}%
\special{pa 2050 2660}%
\special{pa 2055 2688}%
\special{pa 2060 2715}%
\special{pa 2065 2743}%
\special{pa 2070 2770}%
\special{pa 2075 2798}%
\special{pa 2075 2800}%
\special{fp}%
% FUNC 2 0 3 0 Black White  
% 10 400 400 2400 2800 800 2400 2000 2400 800 1200 400 401 2400 2800 0 2 0 1 0 0
% 0.5
\special{pn 8}%
\special{pa 1400 400}%
\special{pa 1400 2800}%
\special{fp}%
% FUNC 2 0 3 0 Black White  
% 10 400 400 2400 2800 800 2400 2000 2400 800 1200 400 2400 2400 2400 0 0 0 1 0 0
% 0.25
\special{pn 8}%
\special{pn 8}%
\special{pa 1100 400}%
\special{pa 1100 408}%
\special{fp}%
\special{pa 1100 445}%
\special{pa 1100 454}%
\special{fp}%
\special{pa 1100 491}%
\special{pa 1100 499}%
\special{fp}%
\special{pa 1100 536}%
\special{pa 1100 544}%
\special{fp}%
\special{pa 1100 582}%
\special{pa 1100 590}%
\special{fp}%
\special{pa 1100 627}%
\special{pa 1100 635}%
\special{fp}%
\special{pa 1100 673}%
\special{pa 1100 681}%
\special{fp}%
\special{pa 1100 718}%
\special{pa 1100 726}%
\special{fp}%
\special{pa 1100 764}%
\special{pa 1100 772}%
\special{fp}%
\special{pa 1100 809}%
\special{pa 1100 817}%
\special{fp}%
\special{pa 1100 855}%
\special{pa 1100 863}%
\special{fp}%
\special{pa 1100 900}%
\special{pa 1100 908}%
\special{fp}%
\special{pa 1100 945}%
\special{pa 1100 954}%
\special{fp}%
\special{pa 1100 991}%
\special{pa 1100 999}%
\special{fp}%
\special{pa 1100 1036}%
\special{pa 1100 1044}%
\special{fp}%
\special{pa 1100 1082}%
\special{pa 1100 1090}%
\special{fp}%
\special{pa 1100 1127}%
\special{pa 1100 1135}%
\special{fp}%
\special{pa 1100 1173}%
\special{pa 1100 1181}%
\special{fp}%
\special{pa 1100 1218}%
\special{pa 1100 1226}%
\special{fp}%
\special{pa 1100 1264}%
\special{pa 1100 1272}%
\special{fp}%
\special{pa 1100 1309}%
\special{pa 1100 1317}%
\special{fp}%
\special{pa 1100 1355}%
\special{pa 1100 1363}%
\special{fp}%
\special{pa 1100 1400}%
\special{pa 1100 1408}%
\special{fp}%
\special{pa 1100 1445}%
\special{pa 1100 1454}%
\special{fp}%
\special{pa 1100 1491}%
\special{pa 1100 1499}%
\special{fp}%
\special{pa 1100 1536}%
\special{pa 1100 1544}%
\special{fp}%
\special{pa 1100 1582}%
\special{pa 1100 1590}%
\special{fp}%
\special{pa 1100 1627}%
\special{pa 1100 1635}%
\special{fp}%
\special{pa 1100 1673}%
\special{pa 1100 1681}%
\special{fp}%
\special{pa 1100 1718}%
\special{pa 1100 1726}%
\special{fp}%
\special{pa 1100 1764}%
\special{pa 1100 1772}%
\special{fp}%
\special{pa 1100 1809}%
\special{pa 1100 1817}%
\special{fp}%
\special{pa 1100 1855}%
\special{pa 1100 1863}%
\special{fp}%
\special{pa 1100 1900}%
\special{pa 1100 1908}%
\special{fp}%
\special{pa 1100 1945}%
\special{pa 1100 1954}%
\special{fp}%
\special{pa 1100 1991}%
\special{pa 1100 1999}%
\special{fp}%
\special{pa 1100 2036}%
\special{pa 1100 2044}%
\special{fp}%
\special{pa 1100 2082}%
\special{pa 1100 2090}%
\special{fp}%
\special{pa 1100 2127}%
\special{pa 1100 2135}%
\special{fp}%
\special{pa 1100 2173}%
\special{pa 1100 2181}%
\special{fp}%
\special{pa 1100 2218}%
\special{pa 1100 2226}%
\special{fp}%
\special{pa 1100 2264}%
\special{pa 1100 2272}%
\special{fp}%
\special{pa 1100 2309}%
\special{pa 1100 2317}%
\special{fp}%
\special{pa 1100 2355}%
\special{pa 1100 2363}%
\special{fp}%
\special{pa 1100 2400}%
\special{pa 1100 2400}%
\special{fp}%
\special{pn 8}%
\special{pa 1100 2409}%
\special{pa 1100 2450}%
\special{fp}%
\special{pa 1100 2459}%
\special{pa 1100 2500}%
\special{fp}%
\special{pa 1100 2509}%
\special{pa 1100 2550}%
\special{fp}%
\special{pa 1100 2559}%
\special{pa 1100 2600}%
\special{fp}%
\special{pa 1100 2609}%
\special{pa 1100 2650}%
\special{fp}%
\special{pa 1100 2659}%
\special{pa 1100 2700}%
\special{fp}%
\special{pa 1100 2709}%
\special{pa 1100 2750}%
\special{fp}%
\special{pa 1100 2759}%
\special{pa 1100 2800}%
\special{fp}%
% STR 2 0 3 0 Black White  
% 4 1540 620 1540 720 5 0 1 0
% $1/2$
\put(15.4000,-7.2000){\makebox(0,0){{\colorbox[named]{White}{$1/2$}}}}%
% STR 2 0 3 0 Black White  
% 4 760 870 760 970 5 0 1 0
% $1/(1+k)^2$
\put(7.6000,-9.7000){\makebox(0,0){{\colorbox[named]{White}{$1/(1+k)^2$}}}}%
\end{picture}}%
}
%  \end{center}

従って$k$の値によってこの最大値は以下のようになる．

\subparagraph{(i)$1/(1+k)^2\le1/2\therefore-1+\sqrt{2}\le k$の時}

$t^2=1/(1+k)^2$の時，つまり$P=C$の時$f(t)$は最大である．

\subparagraph{(ii)$1/2\le1/(1+k)^2\therefore-1+\sqrt{2}\ge k$の時}

$t^2=1/2$つまり$t=1/\sqrt{2}$，$P$の$y$座標が$(k+1)/\sqrt{2}$の時，$f(t)$は最大である．
\casesend

以上から，求める$P$の位置は，
\begin{align}
  \begin{cases}
    |PO|=\dfrac{k+1}{\sqrt{2}}&(k\le-1+\sqrt{2}) \\
    P=C&(-1+\sqrt{2}\le k)
  \end{cases}
\end{align}
である．
\newpage
\end{document}
